\documentclass[../lecture1-introduction.tex]{subfiles}

\begin{document}

\section{Programming}

% -------------------------------------------------------------------

% Fix animation to label, then description
\begin{frame}[fragile]{What is a Computer Program?}
    \begin{description}
        \item[Computer Program] A collection of instructions that performs a specific
        task when executed by a computer. \pause
        \item[Algorithm] A part of a computer program that performs a
        well-defined task. \pause
        \item[Software] A collection of computer programs, libraries and
        related data.
    \end{description}
\end{frame}

% -------------------------------------------------------------------

\begin{frame}[fragile]{Problem Solving}
    Recipe to writing programs: \pause
    \begin{enumerate}
        \item Understand the problem. \pause
        \item Think of a solution. \pause
        \item Describe the solution in as much detail as possible. \newline
              You may use diagrams or plain English to do this. \pause
        \item Translate your solution into a program. \pause
        \item Run your program and see if it works. \pause
        \begin{itemize}
            \item Yes? Hurray! Victory! \pause
            \item No? Go back to 1
        \end{itemize}
    \end{enumerate}
\end{frame}

% -------------------------------------------------------------------

\begin{frame}[fragile]{Example}
    Think like a computer! \newline

    What steps do you need to take to draw a smiley face?
\end{frame}

\begin{frame}[fragile]{Example}
    \begin{center}
        \begin{tikzpicture}
            \path[fill=yellow,draw=yellow!75!red] (0,0) circle (1cm);
            \fill[red] (45:5mm) circle (1mm);
            \fill[red] (135:5mm) circle (1mm);
            \draw[line width=1mm,red] (215:5mm) arc (215:325:5mm);
        \end{tikzpicture}
    \end{center}
\end{frame}

\begin{frame}[fragile]{Example Code Appreciation}
    Lets just take a minute to appreciate what it took to make that smiley face
    \begin{columns}[T,onlytextwidth]
        \column{0.5\textwidth}
            \begin{center}
                \begin{tikzpicture}
                    \path[fill=yellow,draw=yellow!75!red] (0,0) circle (1cm);
                    \fill[red] (45:5mm) circle (1mm);
                    \fill[red] (135:5mm) circle (1mm);
                    \draw[line width=1mm,red] (215:5mm) arc (215:325:5mm);
                \end{tikzpicture}
            \end{center}
        \column{0.5\textwidth}
\begin{cppcode}
\begin{center}
\begin{tikzpicture}
\path[fill=yellow,
        draw=yellow!75!red]
        (0,0) circle (1cm);
\fill[red] (45:5mm) circle (1mm);
\fill[red] (135:5mm) circle (1mm);
\draw[line width=1mm,red] (215:5mm) arc (215:325:5mm);
\end{tikzpicture}
\end{center}
\end{cppcode}
    \end{columns}
    And act like it totally didn't take me like ... ~2 hours to figure out how
    to do it.
\end{frame}

% -------------------------------------------------------------------

\end{document}
